\chapter{Medições em Laboratório}

Na segunda parte da prática, o circuito de filtro passa-faixa com realimentação múltipla (MFB) foi montado em \textit{protoboard}. As medições foram realizadas para caracterizar a resposta em regime permanente senoidal do sistema, comparando os sinais de entrada e saída através de diversas frequências de operação.

\section{Procedimento Experimental}
A etapa experimental consistiu na montagem do circuito ativo e na utilização de instrumentos de bancada para a coleta sistemática de dados. O procedimento seguiu uma sequência lógica de conferência de componentes, alimentação do sistema e medição de sinais com o osciloscópio.

\subsection{Preparação do Circuito}

O circuito foi montado utilizando um amplificador operacional alimentado por uma fonte de tensão simétrica ajustada para $V_{CC} = +10$~V e $V_{EE} = -10$~V. Os componentes utilizados foram:
\begin{itemize}
	\item 1 amplificador operacional (modelo 741 ou equivalente);
	\item Resistores de valores nominais: $R_1 = 470~\Omega$, $R_2 = 470~\Omega$ e $R_3 = 47~\text{k}\Omega$;
	\item 2 capacitores de valor nominal $C = 100$~nF;
	\item Gerador de sinais configurado para uma onda senoidal com amplitude de $100$~mV.
\end{itemize}

\begin{figure}[H]
	\centering
	\caption{Circuito montado em protoboard para realização dos testes.}
	\includegraphics[width=0.7\linewidth]{FIGURAS/circuito_montado.jpg}
	\label{fig:circuito_montado}
	\legend{Fonte: Elaborado pelos autores.}
\end{figure}

\subsection{Resultados Experimentais Coletados}

Durante a execução, utilizou-se o osciloscópio para registrar as amplitudes do canal 1 ($V_{out}$) e do canal 2 ($V_{in}$), além da diferença de fase entre os sinais. Os dados obtidos em bancada estão sumarizados na Tabela \ref{tab:dados_experimentais}.

\begin{table}[H]
	\centering
	\caption{Dados experimentais obtidos em laboratório.}
	\label{tab:dados_experimentais}
	\begin{tabular}{c|c|c|c}
		\hline
		\textbf{Frequência (Hz)} & \textbf{$V_{in}$ (mV)} & \textbf{$V_{out}$ (mV)} & \textbf{Fase ($^\circ$)} \\ \hline
		40,0      & 91,4  & 51,86 & -91,0  \\ [3pt]
		100,0     & 94,9  & 142,3 & -90,0  \\ [3pt]
		200,0     & 97,9  & 320,1 & -93,6  \\ [3pt]
		400,0     & 103,7 & 1593  & -108,8 \\ [3pt]
		480,0     & 89,6  & 1799  & 135,2  \\ [3pt]
		480,2     & 100,6 & 4269  & -153,5 \\ [3pt]
		1101,0    & 91,0  & 364   & 96,5   \\ [3pt]
		5500,0    & 91,7  & 61,9  & 91,5   \\ [3pt]
		10984,0   & 91,5  & 31,9  & 89,7   \\ [3pt]
		\hline
	\end{tabular}
	\legend{Fonte: Dados coletados em laboratório.}
\end{table}

Estes valores experimentais permitem o cálculo da magnitude real da função de transferência através da razão $V_{out}/V_{in}$, possibilitando a comparação com o modelo teórico desenvolvido na etapa anterior.